
\chapter{Conclusion}
\label{ch:conclusion}

\section{The Central Message}

Physical AI controllers routinely make decisions based on observations they
have no reason to trust.  A telemetry packet arrives late, or not at all,
or carrying a stale reading---and the controller evaluates its safety
constraints against a state the plant may never have occupied.  The
resulting Observation-Action Safety Gap is invisible to any evaluation
protocol that checks constraints only on the observed trajectory.

This dissertation has shown, in the battery dispatch domain, that the gap
is real, measurable, and dangerous.  More importantly, it has shown that
the gap is \emph{closable}.  The ORIUS principle---make observation
reliability part of the control state, inflate uncertainty in proportion to
unreliability, and project actions onto the tightened feasible set---eliminates
true-state violations without replacing the upstream optimiser or requiring
a new plant model.

Under degraded observation, safety cannot be taken for granted.  ORIUS makes
observation reliability a first-class quantity in the control loop, not an
afterthought logged for post-hoc analysis.

\section{What the Battery Domain Taught Us}

Grid-scale battery dispatch was chosen as the first proving ground because
it combines hard state constraints, high action-to-state sensitivity, and
network-mediated telemetry.  The experimental campaign yielded several
lessons that extend beyond batteries:

\begin{enumerate}
  \item \textbf{Hidden violations are common.}  Under 20\% dropout, a
    deterministic LP violates true-state SOC bounds at 3.9--4.2\% of
    steps while reporting full feasibility.  The controller believes it
    is safe; the plant disagrees.
  \item \textbf{Conformal prediction is the right uncertainty language.}
    Distribution-free coverage guarantees allow the safety layer to operate
    without a parametric model of telemetry faults.  The CQR intervals
    adapt to the empirical residual distribution, including heavy tails
    induced by sensor drift.
  \item \textbf{Reliability-adaptive inflation is the key mechanism.}
    A conformal interval of fixed width is either too wide when telemetry
    is good (excessive conservatism) or too narrow when telemetry is
    degraded (missed violations).  Scaling width by $1/w_t$ tracks the
    actual risk, concentrating intervention in the episodes that need it.
  \item \textbf{The cost of safety is modest.}  Intervention rates of
    2.8--3.1\% under the canonical fault scenario translate to fewer than
    three corrections per hundred dispatch intervals---a negligible
    operational burden for the elimination of all hidden violations.
\end{enumerate}

\section{Why Observation Reliability Belongs Inside Control}

Classical control treats observation noise as an exogenous disturbance to
be filtered, not a quantity the controller should reason about.  This
separation is justified when the noise model is accurate and stationary.
In modern CPS, neither condition holds: telemetry faults are non-stationary,
packet-level, and correlated with system load.  Treating reliability as
external means the controller cannot tighten its actions when observation
quality degrades, and it cannot relax them when quality is high.

ORIUS brings reliability inside the loop.  The OQE score $w_t$ is updated
at every step from raw telemetry features.  It feeds directly into the
uncertainty inflator, which sets the width of the safe-action region.  The
controller's conservatism is therefore \emph{proportional to its ignorance}:
when it knows little, it acts cautiously; when it knows much, it acts
freely.  This proportionality is the architectural insight of the
dissertation, and it is domain-independent.

\section{What ORIUS Can Become}

The battery result proves the concept.  The adapter architecture
(Chapter~\ref{ch:battery-to-orius}) provides the blueprint for
generalisation.  If the ORIUS community delivers on the roadmap of
Chapter~\ref{ch:roadmap}---a second domain, a public benchmark,
a deployable runtime---the result is a certification layer that any
cyber-physical controller can use, regardless of its internal optimisation
logic.

This is not a small ambition.  It requires that the conformal coverage
guarantee survives domain transfer, that the adapter interface is
expressive enough for nonlinear plants with high-dimensional state, and
that the runtime meets hard real-time constraints on embedded hardware.
Each of these is a non-trivial research or engineering challenge.  But
the structural foundations---domain-abstract scoring, distribution-free
uncertainty, worst-case projection---are in place.

\section{Final Remarks}

This dissertation began with a question: what happens when a controller
cannot trust its own eyes?  The answer, demonstrated rigorously in one
domain and argued structurally for others, is that safety must be
\emph{earned} at every time step by measuring observation quality,
quantifying the resulting uncertainty, and restricting actions to those
that remain safe under the worst case.

The battery domain was the right place to start.  It will not be the
last place ORIUS is needed.
